\subsection{2.11 Symmetrie, Symmetriebruch und die Rekonstruktion der Physik}

Mit der Entwicklung der Projektionstheorie änderte sich die Forschungsstrategie der PST ein weiteres Mal.

Bis zu diesem Zeitpunkt lautete die Fragestellung:

\[
\boxed{
\text{Welche Projektion erzeugt eine vierdimensionale Raumzeit?}
}
\]

Nun entstand eine neue Idee.

Nicht die Raumzeit sollte Ausgangspunkt sein.

Sondern die bekannte Physik selbst.

%%%%%%%%%%%%%%%%%%%%%%%%%%%%%%%%%%%%%%%%%%%%%%%%%%%%%%%%%%%%%%

\section*{Der Perspektivwechsel}

Die bisherige Vorgehensweise war konstruktiv:

\[
\text{OPP}
\longrightarrow
\text{Projektion}
\longrightarrow
\text{Raumzeit}
\longrightarrow
\text{Physik}.
\]

Jetzt wurde erstmals die umgekehrte Richtung betrachtet.

\[
\boxed{
\text{Bekannte Physik}
\longrightarrow
?
\longrightarrow
\text{OPP}.
}
\]

Statt aus dem OPP etwas erzeugen zu wollen,

wurde gefragt:

\begin{quote}
Welche Eigenschaften muss der OPP besitzen, damit unsere bekannte Physik überhaupt entstehen kann?
\end{quote}

Dies entspricht einem klassischen Top-Down-Ansatz.

%%%%%%%%%%%%%%%%%%%%%%%%%%%%%%%%%%%%%%%%%%%%%%%%%%%%%%%%%%%%%%

\section*{Symmetrien als empirische Anker}

Die moderne Physik besitzt eine Vielzahl außerordentlich erfolgreicher Symmetrieprinzipien.

Beispiele sind

\[
SU(3),
\]

\[
SU(2),
\]

\[
U(1),
\]

sowie

\[
SO(1,3).
\]

Diese Gruppen beschreiben

\begin{itemize}

\item

die starke Wechselwirkung,

\item

die elektroschwache Wechselwirkung,

\item

die elektromagnetische Eichsymmetrie,

\item

und die Lorentzsymmetrie der Raumzeit.

\end{itemize}

Die PST muss diese Strukturen nicht voraussetzen,

aber sie sollte erklären können,

warum gerade diese Symmetrien erscheinen.

%%%%%%%%%%%%%%%%%%%%%%%%%%%%%%%%%%%%%%%%%%%%%%%%%%%%%%%%%%%%%%

\section*{Symmetriebruch}

Ebenso wichtig wie die Symmetrien selbst

ist deren spontane Brechung.

Die Natur besitzt keine vollständig symmetrische Gestalt.

Stattdessen entstehen

\begin{itemize}

\item

Massen,

\item

Teilchenfamilien,

\item

Wechselwirkungen,

\item

und beobachtbare Raumzeit

erst nach verschiedenen Symmetriebrüchen.

Damit stellte sich unmittelbar die Frage:

\[
\boxed{
\text{Kann Projektion selbst als Symmetriebruch verstanden werden?}
}
\]

%%%%%%%%%%%%%%%%%%%%%%%%%%%%%%%%%%%%%%%%%%%%%%%%%%%%%%%%%%%%%%

\section*{Der hochdimensionale OPP}

Die bisherigen numerischen Simulationen sprechen dafür,

dass der rohe OPP eine wesentlich höhere effektive Dimension besitzt als die beobachtbare Raumzeit.

Typische Werte lagen zwischen

\[
d_{\mathrm{eff}}
\approx
30
-
45
\]

je nach Punktzahl und Kernelfunktion.

Diese hohe Dimension erscheint zunächst problematisch.

Unter dem Blickwinkel der Symmetrie erhält sie jedoch eine völlig neue Bedeutung.

Ein hochdimensionaler Raum besitzt grundsätzlich wesentlich größere Symmetriegruppen.

Die beobachtete Physik könnte daher als gebrochene Reststruktur einer viel größeren Symmetrie verstanden werden.

%%%%%%%%%%%%%%%%%%%%%%%%%%%%%%%%%%%%%%%%%%%%%%%%%%%%%%%%%%%%%%

\section*{Die Rolle der Projektion}

An dieser Stelle verschmelzen zwei zuvor getrennte Ideen.

Bisher wurde angenommen:

\[
\text{Projektion}
=
\text{Informationsreduktion}.
\]

Nun kommt hinzu:

\[
\boxed{
\text{Projektion}
=
\text{Symmetriebruch}.
}
\]

Eine Projektion reduziert nicht nur Freiheitsgrade.

Sie zerstört gleichzeitig einen Teil der ursprünglichen Symmetrien.

Die beobachtbare Physik wäre dann genau jene Reststruktur,

die unter dieser Projektion erhalten bleibt.

%%%%%%%%%%%%%%%%%%%%%%%%%%%%%%%%%%%%%%%%%%%%%%%%%%%%%%%%%%%%%%

\section*{Eine erste Vermutung}

Während der Diskussion entstand erstmals folgende Hypothese.

Vielleicht besitzt der OPP eine sehr große kontinuierliche Symmetriegruppe

\[
G.
\]

Die Projektion reduziert diese Gruppe auf eine Untergruppe

\[
H.
\]

Formal:

\[
\boxed{
G
\longrightarrow
H.
}
\]

Die beobachtbaren Eichsymmetrien wären genau diese Restgruppe.

%%%%%%%%%%%%%%%%%%%%%%%%%%%%%%%%%%%%%%%%%%%%%%%%%%%%%%%%%%%%%%

\section*{Der Zusammenhang mit Kaluza-Klein}

An dieser Stelle wurde eine bemerkenswerte Parallele zur klassischen Kaluza-Klein-Theorie erkannt.

Dort entstehen Eichfelder dadurch,

dass zusätzliche Dimensionen kompaktifiziert werden.

In der PST wäre eine Kompaktifizierung möglicherweise gar nicht notwendig.

Die Projektion selbst könnte dieselbe mathematische Wirkung besitzen.

Die zusätzlichen Freiheitsgrade verschwinden nicht.

Sie erscheinen lediglich nicht mehr als makroskopische Raumdimensionen.

%%%%%%%%%%%%%%%%%%%%%%%%%%%%%%%%%%%%%%%%%%%%%%%%%%%%%%%%%%%%%%

\section*{Lie-Gruppen im OPP}

Ein weiterer Gedanke entstand aus der Beobachtung,

dass die theoretische Physik ohnehin ständig mit hochdimensionalen Gruppen arbeitet.

Beispiele:

\[
SU(3),
\]

\[
SO(N),
\]

\[
Spin(N),
\]

\[
E_8,
\]

usw.

Dies führte zur Vermutung,

dass die hohe Dimension des OPP gar kein Problem,

sondern möglicherweise sogar eine notwendige Voraussetzung ist.

Die innere Struktur dieser Gruppen könnte bereits vollständig im OPP enthalten sein.

%%%%%%%%%%%%%%%%%%%%%%%%%%%%%%%%%%%%%%%%%%%%%%%%%%%%%%%%%%%%%%

\section*{Fake-5D}

Im Zusammenhang mit den numerischen Resultaten entstand später die Diskussion über die sogenannte

\begin{quote}
„Fake-5D“
\end{quote}

Idee.

Die numerischen Untersuchungen zeigten häufig

\[
d_{\mathrm{eff}}
\approx
5
\]

anstatt exakt

\[
4.
\]

Daraus entstand zunächst die Vermutung,

es könne tatsächlich eine zusätzliche fünfte Dimension existieren.

Später wurde jedoch eine alternative Interpretation bevorzugt.

Nicht fünf reale Raumdimensionen,

sondern

vier makroskopische Projektionsdimensionen

plus

eine effektive Restmode

des Spektrums.

Diese zusätzliche Mode könnte beispielsweise

\begin{itemize}

\item

innere Freiheitsgrade,

\item

Projektionsunschärfe,

\item

oder Quantenfluktuationen

repräsentieren.

%%%%%%%%%%%%%%%%%%%%%%%%%%%%%%%%%%%%%%%%%%%%%%%%%%%%%%%%%%%%%%

\section*{Die zentrale Forschungsfrage}

Aus diesen Überlegungen entstand schließlich eine neue Zielsetzung.

Nicht mehr:

\[
\text{„Wie entsteht Raum?“}
\]

sondern

\[
\boxed{
\text{Welche Symmetrien bleiben unter physikalischen Projektionen erhalten?}
}
\]

Diese Fragestellung verbindet

\begin{itemize}

\item

die numerischen OPP-Simulationen,

\item

die Projektionstheorie,

\item

die Eichsymmetrien des Standardmodells,

\item

sowie die geometrischen Grundlagen der Relativitätstheorie.

%%%%%%%%%%%%%%%%%%%%%%%%%%%%%%%%%%%%%%%%%%%%%%%%%%%%%%%%%%%%%%

\section*{Die Bedeutung für die PST}

Dieser Perspektivwechsel markiert einen weiteren Wendepunkt der Theorie.

Die PST sucht nicht länger direkt nach Raum.

Sie sucht nach derjenigen Projektionsstruktur,

deren erhaltene Symmetrien exakt unserer beobachteten Physik entsprechen.

Raumzeit erscheint damit nicht mehr als primäres Objekt,

sondern als eine Konsequenz gebrochener Symmetrien innerhalb eines wesentlich größeren ontophänomenalen Phasenraums.

%%%%%%%%%%%%%%%%%%%%%%%%%%%%%%%%%%%%%%%%%%%%%%%%%%%%%%%%%%%%%%

\section*{Ausblick}

Die hier entwickelten Überlegungen besitzen bislang überwiegend heuristischen Charakter.

Sie liefern jedoch erstmals einen klaren Fahrplan für zukünftige Arbeiten.

Insbesondere ergeben sich folgende Forschungsziele:

\begin{enumerate}

\item

Mathematische Beschreibung von Projektionen als Symmetriebrechung.

\item

Analyse der erhaltenen Untergruppen.

\item

Verbindung zwischen Projektionsoperatoren und Eichgruppen.

\item

Untersuchung, ob Lorentz-Symmetrie ebenfalls als Restgruppe entsteht.

\item

Numerische Simulation projektionsinduzierter Symmetriebrüche.

\end{enumerate}

Damit endet die erste Entwicklungsphase der PST.

Die Theorie besitzt nun

\begin{itemize}

\item

ein ontologisches Fundament (OP),

\item

einen hochdimensionalen OPP,

\item

eine Projektionstheorie,

\item

erste numerische Evidenz,

\item

und einen klaren physikalischen Forschungsweg.

Die folgenden Bände widmen sich der mathematischen Ausarbeitung dieser Struktur und ihrer Verbindung zur etablierten theoretischen Physik.

